% !TEX program = pdflatex
\documentclass[11pt]{article}
\usepackage{support-notes}
\title{Linear-recurrence zero sets in the Ext of finite modules}
\author{\workingauthor}\date{\coursedate}
\begin{document}\maketitle
\begin{abstract}
Every scalar linear recurrence over a field can be encoded, up to an
initial segment, by Ext between ordinary finite modules over a split
local algebra with radical cube zero. The source is cyclic and has a
totally acyclic rank-one complete resolution; the target has dimension
two. A six-dimensional example over $\mathbb F_p(t)$ gives an
irrational pair Poincar\'e series and arbitrarily long zero intervals.
For the larger class of normal-element dg kernels there is a growth
dichotomy, an eventual periodicity theorem in characteristic zero and
an automaticity theorem in positive characteristic. Results below
are proved.
\end{abstract}

\section{An elementary realisation in ordinary modules}
\begin{thm}\label{recur}
Given $Q\in{\rm GL}_r(k)$, $v\in k^r$, a row $\ell\in (k^r)^*$ and
$c\in k$, set $a_j=c+\ell Q^jv$ and $b_j=\mathbf1_{a_j=0}$.
There is a split local $k$-algebra $L_Q$ of dimension $2r+2$ with
radical cube zero and finite right modules $M_v,N_{c,\ell}$ of
dimensions $r+1,2$ such that for every $j\ge1$,
\[
             \dim_k\Ext^j_{L_Q}(M_v,N_{c,\ell})=b_{j-1}+b_j .
                                                        \tag{1}
\]
The minimal free resolution of $M_v$ extends in both directions to
a totally acyclic complex of free rank-one modules. Thus every scalar
linear recurrence of order $d\ge1$ is represented by (1) in all
sufficiently high degrees with $\dim L_Q\le 2d+2$.
\end{thm}
\begin{proof}
Take basis $1,a,b_1,\ldots,b_r,w_1,\ldots,w_r$. The only nonzero
radical products are
\[
              a b_i=w_i,\qquad b_i a=-\sum_j Q_{ji}w_j .
\]
This table is associative since all triple radical products vanish.
The radical is the span of the nonunit basis, its square is
$W=\langle w_i\rangle$, and its quotient field is $k$.
For a column $s=(s_i)$ write $l(s)=a+\sum_i s_i b_i$. Then
\[
 l(s)l(s')=\sum_j(s'_j-(Qs)_j)w_j,\qquad
 l(s)L_Q=k l(s)\oplus W.
\]
Left multiplication by $l(s)$ on the right regular module has rank
$r+1$, so its kernel has dimension $r+1$. It contains
$l(Qs)L_Q$ by the product formula, hence is exactly that ideal.
Set $M_v=L_Q/l(v)L_Q$. Its free resolution has one copy of $L_Q$
in each homological degree and successive maps to the left given
by left multiplication by $l(v),l(Qv),l(Q^2v),\ldots$.
All entries belong to the radical, so the resolution is minimal.
Since $Q$ is invertible, it continues exactly to negative degrees.
Also $L_Q l(s)=k l(s)\oplus W$, as the $b_i l(s)$ span $W$ by
invertibility of $Q$. Thus the kernel of right multiplication by
$l(s)$ is $L_Q l(Q^{-1}s)$ by the same dimension argument.
Dualizing the doubly infinite complex into $L_Q$ yields exactly
these right multiplications and is exact. This proves total acyclicity
(and hence total reflexivity of its syzygies).

Give $N$ basis $e,f$ with $e a=c f$, $e b_i=\ell_i f$.
The radical kills $f$, and $W$ kills $N$, so this defines a module,
also for the zero functional. In the Hom complex each term is $N$,
and the differential from degree $j$ to $j+1$ sends $e$ to $a_jf$
and $f$ to zero. Its rank is $1-b_j$ and successive differentials
compose to zero. Subtracting the two adjacent ranks from $\dim N=2$
gives (1). In degree zero the dimension is $1+b_0$.

Realize any recurrence by its companion matrix $J$ on $k^d$, an
initial column and a coordinate row. In the Fitting decomposition
for $J$, its nilpotent summand contributes zero once $j\ge d$.
Restricting to the stable image gives the required invertible $Q$
and $c=0$. If this image is zero take $r=1$, $Q=1$, $v=0$.
Only a finite initial segment can change. No factorisation or
separability of the characteristic polynomial is needed.
\end{proof}

This algebra sits in the normal-element class. In the reduced bar-dual
convention with degree-zero products giving degree-two words, the
differentials are $dy_j=xz_j-\sum_i Q_{ji}z_i x$, with $x,z_i,y_j$
all of degree one. Leading words $xz_j$ have no overlaps, so the
dual projects quasi-isomorphically to $A_{1,r,Q}$ by the filtered
word proof in \othernote{normal_localisations}{Bounded profiles and
finite-dimensional points of universal localisations}. Reversing the
bar word convention replaces the presentation by its opposite; the
right-module proof above fixes the side independently. In particular
the residue-field Poincar\'e series of $L_Q$ is
$1/((1-u)(1-r u))$ even in the next example.

\begin{cor}\label{frobenius}
Over $k=\mathbb F_p(t)$ there is a six-dimensional ordinary algebra
as above with a cyclic three-dimensional source and two-dimensional
target whose pair Poincar\'e series over $\mathbb Q$ is
\[
         1+(1+u)\sum_{a\ge0}u^{p^a}.
\]
The groups have unbounded gaps and infinitely many nonzero terms.
In particular eventual disjointness from a single generic divisor
or vanishing through a fixed finite window cannot certify boundedness.
\end{cor}
\begin{proof}
Set $Q=\operatorname{diag}(t,t+1)$, $v=(1,1)^{\rm tr}$,
$\ell=(-1,1)$, $c=-1$. The scalar is
$(t+1)^j-t^j-1$. Its zeros are precisely the positive powers
$p^a$, including $p^0=1$: write $j=p^a m$ with $p\nmid m$ and,
if $m>1$, compare the coefficient of $t^{p^a}$. Also $a_0=-1$.
Formula (1) gives the asserted integer series and $\dim\operatorname{Hom}=1$.
A rational series has an eventual linear recurrence. A sufficiently
long consecutive run of zeros in such a recurrence forces the rest
of the tail to vanish. This contradicts the next power of $p$.
For the geometric failure see the explicit dense two-dimensional
torus orbit in \othernote{orbit_closures}{Orbit closures on
representation spaces}. Here each term of the source's minimal
resolution still has rank one.
\end{proof}

\section{A growth and degree-support theorem for the whole family}
Let $A=A_{n,r,Q}$ for $n,r\ge1$ and $Q\in{\rm GL}_r(k)$ arbitrary.
For $P,E\in\per A$ set $Y=P^\vee$, $G=A/(x^n)$ and
$e_i=\dim H^i\RH_A(P,E)$. The same statements hold for the corresponding
objects in $D_{\fd}$ of its finite-dimensional dg model.
\begin{thm}\label{dichotomy}
If neither $H^*E$ nor $H^*Y$ has eventually bijective $x^n$, then
there is a fixed $b$ with $e_{i+b}\ge\dim G^i$ for all $i\ge0$.
Otherwise the $e_i$ are bounded for $i\gg0$.
In this second case they are eventually periodic in characteristic
zero and form a $p$-automatic sequence up to an initial segment in
characteristic $p$. Consequently in the whole family the set of
nonzero high Ext degrees is eventually periodic in characteristic
zero and $p$-automatic in characteristic $p$. For $nr\ge2$ the
first alternative has exponential growth.
\end{thm}
\begin{proof}
The first estimate is the relative bound in \othernote{two_cuts}{A
lower bound for derived tensor products}. Here $G=(k[x]/x^n)*
k\langle z_1,\ldots,z_r\rangle$ contains $z_1^i$ for every $i$.
If $r\ge2$ its words include the free words in two $z$'s. If $n\ge2$
they include words in the two distinct blocks $z_1,xz_1$; for example
$z_1z_1,xz_1$ give two independently concatenable length-two blocks.
Thus $nr\ge2$ gives an exponential lower bound along a subsequence
(indeed in all degrees by appending $z_1$).

In the second case the normal-element test identifies the tensor with
its localization on high cohomology. The tame side gives a fixed finite
module on the chart, and the other side is the cone of a finite matrix.
In degree $i$ evaluation on the twisted tuple gives a linear map
between vector spaces of fixed finite dimensions. Each $e_i$ is the
sum of a kernel dimension and the preceding cokernel dimension,
up to a fixed shift and choice of cone convention; equivalently for
a tensor cone one has the cokernel in degree $i$ and kernel in degree
$i+1$. Hence the $e_i$ are bounded. Each rank inequality is a finite
Boolean combination of polynomial-minor vanishings in the iterated
tuple. In characteristic zero the return lemma in the orbit-closure
note makes each such set eventually periodic. In characteristic $p$
the linear action on the $nr$ generators makes each evaluated minor
a linear recurrence by the same tensor-power argument as for a full
determinant there. Derksen's theorem makes the zero sets $p$-automatic.
These sets are closed under finite Boolean operations and integer
shifts (finite automata can implement addition with carry). There are
only finitely many rank values. This proves the assertion for the
dimension sequence. In the first case every sufficiently high Ext
degree is nonzero, proving the final assertion about both cases.
\end{proof}

\paragraph{Realisation with genuine two-term dg cohomology.}
The paired zero set can also be realised without an adjacent-index
sum at a doubled index. Start with $a_j=\ell Q^jv$ and take
$A_{2,r+1,Q'}$, $Q'=\operatorname{diag}(Q,1)$. Lift the one-dimensional
right chart tuple with even slot $(v,0)$ and odd slot $(0,\ldots,0,1)$
to the perfect nonnegative tail module $M$. Let $P$ be the cone of
$\sum_i\ell'_i z_i:A[-1]\to A$, $\ell'=(\ell,1)$.
In successive degrees of $M$ the map on Hom from its two free terms
is the scalar $a_j$ in degree $2j$ and $1$ in degree $2j+1$.
The kernel--cokernel sequence gives
$\dim H^{2j}\RH(P,M)=\dim H^{2j+1}\RH(P,M)=\mathbf1_{a_j=0}$
for $j\ge1$. The strict two-term realization of this $A$ has graded
cohomology independent of $Q'$. This provides prescribed recurrence
tails inside the two-term setting of the manuscript.

\paragraph{Priority and limits.}
Parameter-dependent syzygies in the quantum exterior algebra are
classical, including the work of Schulz and
\href{https://arxiv.org/abs/1511.01418}{Erdmann}. Unbounded delays
before eventual Ext vanishing occur in
\href{https://arxiv.org/abs/math/0306001}{Jorgensen--\c Sega}.
\href{https://arxiv.org/abs/1706.06470}{Piontkovski} realises
linear-orbit return times in Hilbert functions, including irrational
series in positive characteristic. The arithmetic example above is
Lech's, discussed by \href{https://arxiv.org/abs/math/0510583}{Derksen}.
The additional assertion here is the ordinary-module formula (1)
with the stated algebra size, cyclic totally reflexive source and
two-dimensional target, together with the classification consequences
for the dg family. The arithmetic is an input. We make no priority
claim for irrational homological series in general or for arbitrary
sets of integers. Eventual boundedness tests infinitely many zeros
rather than existence of a single zero in a recurrence.
\end{document}
